\documentclass[11pt,a4paper]{article}

% ============================================================
%  QIMMA -- COURS : Équations différentielles
%  2ème Bac Sciences Mathématiques (A et B)
%  Standard exhaustif : définitions formelles + démonstrations
%  Basé sur templates/qimma-template.tex (charte Qimma)
% ============================================================

\usepackage[french]{babel}
\usepackage[utf8]{inputenc}
\usepackage[T1]{fontenc}
\usepackage{lmodern}
\usepackage[a4paper,margin=2cm,top=2.3cm,bottom=2.6cm]{geometry}
\usepackage{amsmath,amssymb}
\usepackage{xcolor}
\usepackage{graphicx}
\usepackage{tikz}
\usetikzlibrary{shadows,calc}
\usepackage[most]{tcolorbox}
\usepackage{titlesec}
\usepackage{fancyhdr}
\usepackage{enumitem}
\usepackage{pifont}
\usepackage{array}
\usepackage{colortbl}
\usepackage{hyperref}

% ------------------- CHARTE QIMMA -------------------
\definecolor{qteal}{HTML}{0d9488}
\definecolor{qtealdark}{HTML}{134e4a}
\definecolor{qamber}{HTML}{fbbf24}
\definecolor{qambertext}{HTML}{92400e}
\definecolor{ink}{HTML}{1f2937}
\definecolor{inkgray}{HTML}{6b7280}
\definecolor{tealpale}{HTML}{e6f5f3}
\colorlet{colDef}{qteal}
\definecolor{colProp}{HTML}{0f766e}
\definecolor{colEx}{HTML}{b45309}
\definecolor{colRem}{HTML}{9d174d}
\color{ink}

\graphicspath{{../../../../assets/}}
\newcommand{\logofile}{../../../../assets/qimma-logo.pdf}

% ------------------- METADONNEES -------------------
\newcommand{\matiere}{Mathématiques}
\newcommand{\niveau}{2\up{ème} Bac -- Sciences Mathématiques}
\newcommand{\chapitretitre}{Équations différentielles}

% ------------------- EN-TETE / PIED DE PAGE -------------------
\pagestyle{fancy}
\fancyhf{}
\renewcommand{\headrulewidth}{0pt}
\renewcommand{\footrulewidth}{0.5pt}
\fancyfoot[L]{\raisebox{-0.15cm}{\includegraphics[height=0.35cm]{\logofile}}~\small\sffamily\color{inkgray}Qimma}
\fancyfoot[C]{\small\sffamily\color{inkgray}\matiere\ -- \niveau}
\fancyfoot[R]{\small\sffamily\color{qtealdark}\thepage}

% ------------------- TITRES DE SECTION -------------------
\titleformat{\section}
  {\normalfont\sffamily\Large\bfseries\color{qtealdark}}
  {\tikz[baseline=-3.2pt]{\node[circle,fill=qteal,text=white,inner sep=0pt,minimum size=8mm]{\sffamily\bfseries\thesection};}}
  {10pt}{}
  [\vspace{-2mm}{\color{qamber}\titlerule[1.6pt]}\vspace{2mm}]
\titlespacing*{\section}{0pt}{16pt}{10pt}

\titleformat{\subsection}
  {\normalfont\sffamily\large\bfseries\color{qtealdark}}
  {\color{qteal}\ding{212}}{6pt}{}
\titlespacing*{\subsection}{0pt}{12pt}{6pt}

% ------------------- BOITES PEDAGOGIQUES -------------------
\newtcolorbox{fiche}[3][]{
  enhanced, breakable,
  colback=white, colframe=#2,
  boxrule=1pt, arc=2mm,
  top=7mm, left=4mm, right=4mm, bottom=3mm,
  drop shadow={black!25, shadow xshift=1mm, shadow yshift=-1mm},
  attach boxed title to top left={xshift=4mm, yshift=-3mm, yshifttext=-1mm},
  boxed title style={colback=#2, arc=1.5mm, boxrule=0pt, left=2mm, right=2mm},
  coltitle=white, fonttitle=\sffamily\bfseries,
  title={#3}, #1
}
\newenvironment{definition}[1][Définition]
  {\begin{fiche}{colDef}{\ding{110}~~#1}}{\end{fiche}}
\newenvironment{propriete}[1][Propriété]
  {\begin{fiche}{colProp}{\ding{72}~~#1}}{\end{fiche}}
\newenvironment{exemple}[1][Exemple]
  {\begin{fiche}{colEx}{\ding{217}~~#1}}{\end{fiche}}
\newenvironment{remarque}[1][Remarque]
  {\begin{fiche}{colRem}{\ding{43}~~#1}}{\end{fiche}}

% Démonstration (hors boîte, style discret)
\newenvironment{demo}
  {\par\smallskip\noindent{\sffamily\bfseries\color{qtealdark}Démonstration.}\ \itshape}
  {\ \normalfont\hfill$\blacksquare$\par\medskip}

\hypersetup{colorlinks=true, linkcolor=qtealdark, urlcolor=qtealdark}

% ============================================================
\begin{document}

% ------------------- BANNIERE DE TITRE -------------------
\begin{tcolorbox}[
  enhanced, boxrule=0pt, arc=4mm,
  interior style={left color=qtealdark, right color=qteal},
  top=8mm, bottom=8mm, left=10mm, right=32mm,
  drop shadow={black!35, shadow xshift=1.5mm, shadow yshift=-1.5mm},
  before skip=0pt, after skip=9mm,
  overlay={
    \node[anchor=north west] at ([xshift=6mm,yshift=-6mm]frame.north west)
      {\includegraphics[height=1.25cm]{\logofile}};
    \node[anchor=north east, fill=qamber, text=qambertext, rounded corners=2pt,
          inner sep=4pt, font=\sffamily\bfseries\small] at ([xshift=-6mm,yshift=-6mm]frame.north east) {COURS};
  }
]
\hspace{1.6cm}\centering
{\sffamily\bfseries\color{white}\fontsize{23}{27}\selectfont \chapitretitre}\\[3mm]
{\sffamily\color{white!90}\large \matiere\ \textbullet\ \niveau}
\end{tcolorbox}

% ------------------- OBJECTIFS -------------------
\begin{tcolorbox}[
  enhanced, colback=tealpale, colframe=qteal, boxrule=0.8pt, arc=2mm,
  left=5mm, right=5mm, top=3mm, bottom=3mm,
  title={\sffamily\bfseries\color{qtealdark}\ding{51}~~Objectifs du chapitre},
  coltitle=qtealdark, colbacktitle=tealpale, boxrule=0.8pt,
  attach title to upper={\par\vspace{1mm}}
]
\begin{itemize}[leftmargin=6mm, label=\ding{212}, itemsep=1mm]
  \item Résoudre et démontrer la solution générale de $y'=ay$ et $y'=ay+b$ ($a\ne0$).
  \item Déterminer la solution vérifiant une condition initiale, et vérifier une solution par substitution.
  \item Résoudre l'équation du second ordre $y''+\omega^2y=0$ avec deux conditions initiales.
  \item Reconnaître ces équations dans des situations de modélisation issues de la physique et de la biologie.
  \item Traiter les exercices type examen national combinant équation différentielle et étude de fonction.
\end{itemize}
\end{tcolorbox}

\vspace{4mm}

% ------------------- SECTION 1 -------------------
\section{L'équation $y'=ay$}

\begin{definition}
Une \textbf{équation différentielle} est une équation dont l'inconnue est une \textbf{fonction} $y$, faisant intervenir $y$ et ses dérivées. \textbf{Résoudre} l'équation sur un intervalle, c'est trouver \textbf{toutes} les fonctions dérivables qui la vérifient sur cet intervalle.
\end{definition}

\begin{propriete}[Solutions de $y'=ay$]
Soit $a\in\mathbb{R}^*$. Les solutions sur $\mathbb{R}$ de $y'=ay$ sont exactement les fonctions
\[
  y(x) = Ce^{ax}, \qquad C\in\mathbb{R}.
\]
\end{propriete}

\begin{demo}
\emph{Les fonctions $Ce^{ax}$ sont solutions} : si $y(x)=Ce^{ax}$, alors $y'(x)=Cae^{ax}=ay(x)$. $\checkmark$

\smallskip
\emph{Ce sont les seules} : soit $y$ une solution de $y'=ay$. Posons $z(x)=y(x)e^{-ax}$. Alors
\[
  z'(x) = y'(x)e^{-ax}+y(x)\times(-a)e^{-ax} = e^{-ax}\bigl(y'(x)-ay(x)\bigr) = e^{-ax}\times0 = 0.
\]
Donc $z$ est \textbf{constante} sur $\mathbb{R}$ (intervalle) : $z(x)=C$, soit $y(x)e^{-ax}=C$, c'est-à-dire $y(x)=Ce^{ax}$.
\end{demo}

\begin{propriete}[Unicité avec condition initiale]
Pour tout $(x_0,y_0)$, il existe une \textbf{unique} solution de $y'=ay$ vérifiant $y(x_0)=y_0$.
\end{propriete}

\begin{exemple}[Résolution avec condition initiale]
Résolvons $y'=-2y$ avec $y(0)=5$.

\smallskip
Solutions générales : $y(x)=Ce^{-2x}$. Condition $y(0)=5$ : $C=5$.

\smallskip
\textbf{Conclusion} : $y(x)=5e^{-2x}$ -- modèle de \textbf{décroissance exponentielle} (retrouvé en physique : décroissance radioactive $N(t)=N_0e^{-\lambda t}$).
\end{exemple}

% ------------------- SECTION 2 -------------------
\section{L'équation $y'=ay+b$}

\begin{propriete}[Solutions de $y'=ay+b$]
Soient $a\in\mathbb{R}^*$, $b\in\mathbb{R}$. Les solutions sur $\mathbb{R}$ de $y'=ay+b$ sont
\[
  y(x) = Ce^{ax}-\frac ba, \qquad C\in\mathbb{R}.
\]
\end{propriete}

\begin{demo}
Posons $y_p=-\frac ba$, la \textbf{solution particulière constante} (vérification : $y_p'=0$ et $ay_p+b=a\left(-\frac ba\right)+b=0$, donc $y_p'=ay_p+b$ est vérifiée). Soit $y$ une solution quelconque de $y'=ay+b$, et posons $z=y-y_p$. Alors
\[
  z' = y'-y_p' = (ay+b)-0 = ay+b = a(z+y_p)+b = az+ay_p+b = az+0 = az
\]
(en utilisant $ay_p+b=0$). Donc $z$ vérifie $z'=az$ : d'après la section précédente, $z(x)=Ce^{ax}$, soit $y(x)=Ce^{ax}+y_p=Ce^{ax}-\frac ba$.
\end{demo}

\begin{remarque}[Méthode en 3 étapes]
\begin{enumerate}[leftmargin=7mm, itemsep=0.5mm]
  \item Chercher la \textbf{solution constante} : $0=a\ell+b$, d'où $\ell=-\frac ba$.
  \item Écrire la solution générale : $y(x)=Ce^{ax}+\ell$.
  \item Déterminer $C$ avec la condition initiale.
\end{enumerate}
Parallèle avec les suites arithmético-géométriques $u_{n+1}=au_n+b$ : même logique de point fixe + partie exponentielle/géométrique.
\end{remarque}

\begin{exemple}[Résolution complète type bac]
Résolvons $y'=3y-6$ avec $y(0)=1$.

\smallskip
\textbf{Solution constante} : $0=3\ell-6\iff\ell=2$.

\smallskip
\textbf{Solution générale} : $y(x)=Ce^{3x}+2$.

\smallskip
\textbf{Condition initiale} : $y(0)=C+2=1\iff C=-1$.

\smallskip
\textbf{Conclusion} : $y(x)=2-e^{3x}$. \emph{Vérification} : $y'=-3e^{3x}$ et $3y-6=6-3e^{3x}-6=-3e^{3x}$. $\checkmark$
\end{exemple}

% ------------------- SECTION 3 -------------------
\section{L'équation $y''+\omega^2y=0$}

\begin{propriete}[Solutions de l'oscillateur harmonique]
Soit $\omega\in\mathbb{R}^*$. Les solutions sur $\mathbb{R}$ de $y''+\omega^2y=0$ sont
\[
  y(x) = A\cos(\omega x)+B\sin(\omega x), \qquad A,B\in\mathbb{R}.
\]
Forme équivalente : $y(x)=R\cos(\omega x-\varphi)$ avec $R\geq0$, $\varphi\in\mathbb{R}$ (amplitude et phase). Une solution unique est déterminée par \textbf{deux} conditions initiales, par exemple $y(0)$ et $y'(0)$.
\end{propriete}

\begin{demo}[Vérification que ces fonctions sont solutions]
Si $y(x)=A\cos(\omega x)+B\sin(\omega x)$ : $y'(x)=-A\omega\sin(\omega x)+B\omega\cos(\omega x)$ et $y''(x)=-A\omega^2\cos(\omega x)-B\omega^2\sin(\omega x)=-\omega^2y(x)$. Donc $y''+\omega^2y=0$. $\checkmark$ (L'unicité, c'est-à-dire que ce sont \emph{toutes} les solutions, est admise à ce niveau.)
\end{demo}

\begin{exemple}[Conditions initiales doubles]
Résolvons $y''+4y=0$ avec $y(0)=1$, $y'(0)=2$.

\smallskip
Ici $\omega=2$ : solution générale $y(x)=A\cos2x+B\sin2x$.

\smallskip
$y(0)=A=1$. Comme $y'(x)=-2A\sin2x+2B\cos2x$, on a $y'(0)=2B=2\iff B=1$.

\smallskip
\textbf{Conclusion} : $y(x)=\cos2x+\sin2x$, que l'on peut écrire $\sqrt2\cos\left(2x-\frac\pi4\right)$.
\end{exemple}

\begin{remarque}[Lien avec la physique]
Cette équation modélise les \textbf{oscillateurs harmoniques} : pendule élastique ($\ddot x+\frac km x=0$), circuit LC ($\ddot q+\frac1{LC}q=0$)\ldots La pulsation $\omega$ donne la période $T=\frac{2\pi}{\omega}$.
\end{remarque}

% ------------------- SECTION 4 -------------------
\section{Exercices corrigés -- niveau examen national SM}

\begin{exemple}[Exercice 1 -- résolution avec vérification]
Vérifier que $y(x)=2e^{-x}+3$ est solution de $y'+y=3$, puis retrouver cette solution par la méthode générale.

\smallskip
\textbf{Corrigé.} \emph{Vérification} : $y'(x)=-2e^{-x}$, donc $y'+y=-2e^{-x}+2e^{-x}+3=3$. $\checkmark$

\smallskip
\emph{Méthode générale} : l'équation s'écrit $y'=-y+3$, soit $a=-1$, $b=3$. Solution constante : $\ell=\frac{-b}{a}=\frac{-3}{-1}=3$. Solution générale : $y(x)=Ce^{-x}+3$. Si $y(0)=5$ (donnée par $y(0)=2+3=5$ dans l'exemple), $C+3=5\iff C=2$, retrouvant $y(x)=2e^{-x}+3$.
\end{exemple}

\begin{exemple}[Exercice 2 -- modélisation : refroidissement (loi de Newton)]
Un objet à température $\theta(t)$ (en $°C$) refroidit dans une pièce à $20°C$ selon la loi $\theta'(t)=-k(\theta(t)-20)$, où $k>0$ est une constante. Sachant $\theta(0)=90$, exprimer $\theta(t)$ en fonction de $t$ et $k$.

\smallskip
\textbf{Corrigé.} On développe : $\theta'=-k\theta+20k$, soit $a=-k$, $b=20k$. Solution constante : $\ell=\frac{-20k}{-k}=20$ (cohérent avec la température ambiante). Solution générale : $\theta(t)=Ce^{-kt}+20$.

\smallskip
$\theta(0)=C+20=90\iff C=70$.

\smallskip
\textbf{Conclusion} : $\theta(t)=70e^{-kt}+20$. Quand $t\to+\infty$, $\theta(t)\to20$ : l'objet se refroidit progressivement jusqu'à la température ambiante.
\end{exemple}

\begin{exemple}[Exercice 3 -- fonction définie par une équation différentielle, étude complète]
Soit $f$ la solution de $y'=2y-4$ telle que $f(0)=1$. Exprimer $f(x)$, puis étudier ses variations et sa limite en $+\infty$ et en $-\infty$.

\smallskip
\textbf{Corrigé.} $a=2$, $b=-4$ : $\ell=\frac{4}{2}=2$. Solution générale $y(x)=Ce^{2x}+2$. $f(0)=C+2=1\iff C=-1$.

\smallskip
\textbf{Conclusion} : $f(x)=2-e^{2x}$.

\smallskip
\emph{Variations} : $f'(x)=-2e^{2x}<0$ pour tout $x$ : $f$ est strictement décroissante sur $\mathbb{R}$.

\smallskip
\emph{Limites} : en $+\infty$, $e^{2x}\to+\infty$, donc $f(x)\to-\infty$. En $-\infty$, $e^{2x}\to0$, donc $f(x)\to2$ (asymptote horizontale $y=2$).
\end{exemple}

\begin{exemple}[Exercice 4 -- second ordre avec interprétation]
Un point matériel oscille selon $x''(t)+9x(t)=0$, avec $x(0)=0$ et une vitesse initiale $x'(0)=6$ (m/s). Déterminer $x(t)$, la période, et l'amplitude du mouvement.

\smallskip
\textbf{Corrigé.} $\omega^2=9\Rightarrow\omega=3$. Solution générale $x(t)=A\cos3t+B\sin3t$.

\smallskip
$x(0)=A=0$. $x'(t)=-3A\sin3t+3B\cos3t$, donc $x'(0)=3B=6\iff B=2$.

\smallskip
\textbf{Conclusion} : $x(t)=2\sin3t$. \textbf{Période} : $T=\frac{2\pi}{3}$ s. \textbf{Amplitude} : $R=\sqrt{A^2+B^2}=2$ m.
\end{exemple}

% ------------------- SECTION 5 -------------------
\section{Méthodes et pièges : la boîte à outils de l'examen}

\subsection{Ce qu'on te demande, ce que tu fais}

\begin{center}
\renewcommand{\arraystretch}{1.45}
\sffamily\small
\begin{tabular}{>{\raggedright\arraybackslash}p{7.2cm} >{\raggedright\arraybackslash}p{8cm}}
\rowcolor{qtealdark} \color{white}\bfseries Ce qu'on te demande & \color{white}\bfseries Ton réflexe \\
\rowcolor{tealpale} Résoudre $y'=ay$ & Solution : $y(x)=Ce^{ax}$. \\
Résoudre $y'=ay+b$ & Solution constante $\ell=-\frac ba$, puis $y(x)=Ce^{ax}+\ell$, puis déterminer $C$. \\
\rowcolor{tealpale} Résoudre $y''+\omega^2y=0$ & $y(x)=A\cos(\omega x)+B\sin(\omega x)$, deux conditions initiales pour $A$ et $B$. \\
Vérifier qu'une fonction est solution & Substituer $y$, $y'$ (et $y''$ si second ordre) dans l'équation, vérifier l'égalité. \\
\rowcolor{tealpale} Modélisation (refroidissement, radioactivité) & Identifier $a$ et $b$ dans l'équation donnée, appliquer la méthode standard. \\
Étudier la fonction solution obtenue & Une fois $y(x)$ explicité, dérivée, signe, limites classiques. \\
\end{tabular}
\end{center}

\subsection{Les pièges qui coûtent des points}

\begin{remarque}[Erreurs relevées chaque année par les correcteurs]
\begin{itemize}[leftmargin=6mm, itemsep=0.5mm]
  \item Oublier le signe « $-$ » dans $\ell=-\frac ba$ (erreur de signe fréquente).
  \item Confondre l'équation $y'=ay+b$ avec sa forme « à l'envers » $y'-ay=b$ : bien identifier $a$ et $b$ en isolant $y'$.
  \item Pour le second ordre : n'utiliser qu'une seule condition initiale (il en faut \textbf{deux}, typiquement $y(0)$ et $y'(0)$).
  \item Oublier de vérifier une solution trouvée en la réinjectant dans l'équation de départ.
  \item Confondre pulsation $\omega$ et période $T$ : $T=\frac{2\pi}{\omega}$ (pas l'inverse).
  \item Dans une modélisation physique, mal identifier la constante $\ell$ avec la grandeur physique attendue (ex. température ambiante, valeur d'équilibre).
\end{itemize}
\end{remarque}

% ------------------- RESUME -------------------
\vspace{4mm}
\begin{tcolorbox}[
  enhanced, boxrule=0pt, arc=2mm,
  interior style={left color=qtealdark, right color=qteal},
  left=6mm, right=6mm, top=4mm, bottom=4mm,
  fontupper=\color{white}\sffamily,
  title={\sffamily\bfseries\color{white}\ding{72}~~À retenir},
  colbacktitle=qtealdark, coltitle=white, boxrule=0pt
]
\begin{itemize}[leftmargin=6mm, label={\color{qamber}\ding{212}}, itemsep=1mm]
  \item $y'=ay$ \ $\Rightarrow$ \ $y=Ce^{ax}$ (démontré via $z=ye^{-ax}$ constante).
  \item $y'=ay+b$ \ $\Rightarrow$ \ $y=Ce^{ax}-\frac ba$ (solution d'équilibre $+$ homogène).
  \item $y''+\omega^2y=0$ \ $\Rightarrow$ \ $y=A\cos\omega x+B\sin\omega x=R\cos(\omega x-\varphi)$.
  \item Une condition initiale fixe $C$ (1\up{er} ordre) ; il en faut \textbf{deux} pour le 2\up{nd} ordre.
  \item Toujours vérifier la solution en la réinjectant dans l'équation de départ.
\end{itemize}
\end{tcolorbox}

\end{document}
